% Rapport d'activités - Candidature à la qualification CNU 27 Informatique
% Lucas Gautheron
\documentclass[11pt,a4paper]{article}

\usepackage[utf8]{inputenc}
\usepackage[T1]{fontenc}
\usepackage[francais]{babel}
\usepackage[top=2.5cm, bottom=2.5cm, left=2.5cm, right=2.5cm]{geometry}
\usepackage{hyperref}
\usepackage{xcolor}
\usepackage{titlesec}
\usepackage{enumitem}
\usepackage{booktabs}
\usepackage{longtable}
\usepackage{array}
\usepackage{fancyhdr}
\usepackage{lastpage}
\usepackage{tabularx}
\usepackage{multirow}

% Colors
\definecolor{DarkGreen}{RGB}{1,125,32}

\definecolor{moderncvgreen}{rgb}{0.35,0.70,0.30}

% Hyperref setup
\hypersetup{
colorlinks=true,
linkcolor=DarkGreen,
urlcolor=DarkGreen,
,citecolor=moderncvgreen
,citebordercolor=moderncvgreen
,urlbordercolor=moderncvgreen
}

% Section formatting
\titleformat{\section}{\Large\bfseries\color{moderncvgreen}}{\thesection.}{0.5em}{}
\titleformat{\subsection}{\large\bfseries\color{moderncvgreen}}{\thesubsection.}{0.5em}{}
\titleformat{\subsubsection}{\normalsize\bfseries\itshape}{\thesubsubsection.}{0.5em}{}

% Header/Footer
\pagestyle{fancy}
\fancyhf{}
\rhead{\textit{\small Rapport d'activités -- Lucas Gautheron -- CNU 27}}
\cfoot{\thepage/\pageref{LastPage}}
\renewcommand{\headrulewidth}{0.4pt}

\begin{document}

%===============================================================================
% TITLE PAGE
%===============================================================================

\begin{center}
{\LARGE\bfseries RAPPORT D'ACTIVITÉS}\\[0.5cm]
{\large Candidature à la qualification aux fonctions de}\\[0.2cm]
{\Large\bfseries Maître de Conférences}\\[0.3cm]
{\large Section 27 -- Informatique}\\[1cm]
\rule{\textwidth}{1pt}\\[0.5cm]
{\Huge\bfseries Lucas GAUTHERON}\\[0.5cm]
\rule{\textwidth}{1pt}\\[1cm]
{\large Décembre 2025}
\end{center}

\vspace{1cm}

\tableofcontents

\newpage

%===============================================================================
% 1. IDENTITÉ
%===============================================================================

\section{Identité}

\begin{tabular}{@{}ll}
\textbf{Nom, Prénom :} & GAUTHERON Lucas \\
\textbf{Date de naissance :} & 29 janvier 1995 \\
\textbf{Nationalité :} & Française \\
\textbf{Courriel :} & \href{mailto:lucas.gautheron@gmail.com}{lucas.gautheron@gmail.com} \\
\textbf{Téléphone :} & +33 6 79 23 86 47 \\
\textbf{Site web :} & \url{https://lucasgautheron.github.io} \\
\textbf{ORCID :} & \href{https://orcid.org/0000-0002-3776-3731}{0000-0002-3776-3731} \\
\end{tabular}

\vspace{0.5cm}

Ma recherche consiste à développer des méthodes computationelles pour l'étude des phénomènes de cognition collective et d'évolution culturelle. 
Ces méthodes comprennent la modélisation formelle de comportement collectifs, ainsi que l'analyse de données, qu'elles soient observationnelles ou issues d'expériences collectives à grande échelle. Bien que ma production soit très interdisciplinaire, l'élaboration et l'exploitation de techniques computationnelles innovantes au service des sciences humaines et sociale est le coeur de mon activité, en collaboration étroite avec des philosophes, des linguistes, ou plus largement des chercheurs en sciences cognitives et anthropologie.



\

\noindent\textbf{Domaines thématiques (nomenclature)}
\begin{itemize}[nosep]
    \item Intelligence artificielle: Intelligence artificielle distribuée, systèmes multi-agents, modélisation cognitive; Théorie de la décision, théorie du choix social; Science des données; Traitement automatique des langues et de la parole
\end{itemize}

\ 

\noindent\textbf{Mots-clés :} modélisation, inférence bayésienne, modèles à variables latentes, problèmes inverses, simulations par agent, analyse de réseaux, traitement automatique des langues, sciences sociales computationnelles.

%===============================================================================
% 2. PARCOURS UNIVERSITAIRE
%===============================================================================

\section{Parcours universitaire}

\subsection{Diplômes obtenus}

\begin{description}[style=nextline, leftmargin=0.5cm]

\item[\textbf{2022--2025} : Doctorat en Sociologie Computationnelle] \hfill \textit{Summa cum laude}\\
\textbf{Établissement :} Interdisciplinary Centre for Science and Technology Studies (IZWT), Bergische Universität Wuppertal, Allemagne\\
\textbf{Titre :} \textit{``Social dilemmas in science: a computational approach to collective action problems and tradeoffs in collective cognition''}\\
\textbf{Directeurs :} Thomas Heinze (Professeur, Sociologie) et Radin Dardashti (Professeur, Philosophie des sciences)\\
\textbf{Jury :} Mark Lutter, Hugo Mercier, Camille Roth, Dunja Šešelja\\
\textbf{Financement :} Research Training Group ``Transformations of Science and Technology since 1800'' (DFG GRK 2696)\\
\textbf{Contribution computationnelle:} Développement et application de méthodes computationnelles (apprentissage automatique, modélisation statistique bayésienne, analyse de réseaux, traitement automatique des langues) pour l'étude des comportements collectifs.

\item[\textbf{2021--2022} : Master 2 Histoire et Philosophie des Sciences] \hfill \textit{Rang : 1\textsuperscript{er}}\\
\textbf{Établissement :} Université Paris-Cité\\
\textbf{Mémoire :} \textit{``Too Beautiful to be False, or Too Beautiful to be True? Searching for supersymmetry at the Large Hadron Collider''}\\
\textbf{Directeurs :} Olivier Darrigol (DR CNRS) et Elisa Omodei (Professeure, Central European University)

\item[\textbf{2014--2018} : Master 1 Physique Théorique et Expérimentale]

\textbf{Établissement :} École Normale Supérieure Paris-Saclay (anciennement ENS Cachan)\\
\textbf{Statut :} Normalien sur concours

\end{description}

%===============================================================================
% 3. EXPÉRIENCES PROFESSIONNELLES
%===============================================================================

\section{Expériences professionnelles}

\begin{description}[style=nextline, leftmargin=0.5cm]

\item[\textbf{Août 2025 -- Août 2027} : Chercheur postdoctoral]
\textbf{Établissement :} Department of Philosophy, Evolution, Science and Society Group, University of Missouri (États-Unis)\\
\textbf{Statut :} CDD postdoctoral\\
\textbf{Activités :} Intelligence et la créativité collectives; développement d'expériences collectives à grande échelle; modélisation et analyse de données comportementales.

\item[\textbf{Mars 2023 -- Décembre 2024}; \textbf{Septembre 2020 -- Novembre 2021}:\\ Ingénieur de recherche en linguistique computationnelle (temps partiel)]
\textbf{Établissement :} Laboratoire de Sciences Cognitives et Psycholinguistique (LSCP), Département d'Études Cognitives, École Normale Supérieure, Paris\\
\textbf{Statut :} CDD\\
\textbf{Activités :} Création d'un package Python pour la gestion, le stockage et l'analyse de grands corpus d'enregistrements audio (gère un volume de 87,000+ heures d'audio parmi 2,700+ enfants à l'heure actuelle); analyse statistique bayésienne pour l'étude de la relation entre l'expérience et la production vocales des enfants

\item[\textbf{Octobre 2016 -- Janvier 2017} : Stage de recherche en astrophysique]
\textbf{Établissement :} Laboratoire Univers et Théories (LUTH), CNRS, Observatoire de Paris-Meudon\\
\textbf{Activités :} Implémentation de calculs de taux d'interaction faible dans un code de simulation d'effondrement de supernovae (Fortran, C++).

\item[\textbf{Mai 2016 -- Juillet 2016} : Stage de recherche en physique des particules]
\textbf{Établissement :} Laboratoire de Physique Nucléaire et des Hautes Énergies (LPNHE), CNRS, Paris\\
\textbf{Activités :} Analyse diphoton et phénoménologie pour l'expérience ATLAS (Pythia, ROOT, RooFit, C++, Python).

\item[\textbf{Mai 2015 -- Juillet 2015} : Stage de recherche en physique des particules]
\textbf{Établissement :} Laboratoire d'Annecy-le-Vieux de Physique des Particules (LAPP), CNRS\\
\textbf{Activités :} Analyse de données et simulations Monte-Carlo pour l'expérience ATLAS (ROOT, C++).

\end{description}

%===============================================================================
% 4. ACTIVITÉS D'ENSEIGNEMENT
%===============================================================================

\section{Activités d'enseignement}

\subsection{Synthèse des enseignements réalisés}

Le tableau ci-dessous résume l'ensemble de mes activités pédagogiques en informatique et domaines connexes. Bien que mon contrat doctoral financé par le DFG ne me donnait en principe pas l'autorisation d'enseigner, j'ai activement organisé des opportunités de transmettre des notions et outils computationnels auprès de publics d'étudiants et chercheurs non spécialistes.

\vspace{0.3cm}

\small
\noindent\begin{tabularx}{\textwidth}{|p{1.6cm}|p{1cm}|p{2.2cm}|p{1.5cm}|p{0.8cm}|p{2.8cm}|p{1cm}|X|}
\hline
\textbf{Statut} & \textbf{Année} & \textbf{Établissement} & \textbf{Public} & \textbf{Niv.} & \textbf{Matière} & \textbf{Vol. (h éq. TD)} & \textbf{Nature} \\
\hline
Chargé de TD & 2025 & EHESS & Master & M1-M2 & Sciences sociales computationnelles & 16 & Cours/TD \\
\hline
Intervenant & Juil. 2025 & Univ. Wuppertal & Doctorants & D & Analyse de réseaux pour SHS & 2 & Tutoriel \\
\hline
Intervenant & Jan. 2025 & Ruhr Univ. Bochum & Master et doctorants & ~ & Problèmes inverses et modèles agents & 2 & Cours invité \\
\hline
Intervenant & Juin 2024 & UCLA & Doctorants et chercheurs & ~ & Logiciel LongFoRMer & 2 & Tutoriel \\
\hline
\end{tabularx}
\normalsize

\vspace{0.3cm}

\textbf{Volume horaire total : environ 22 heures équivalent TD}

\subsection{Description détaillée des enseignements}

\subsubsection{Sciences sociales computationnelles (EHESS, 2025)}

\textbf{Responsable du cours :} Camille Roth (EHESS)\\
\textbf{Public :} Étudiants de master en sciences sociales\\
\textbf{Rôle :} Co-animateur 

Ce séminaire (\url{https://ssc-ehess.github.io/}) forme les chercheurs en sciences sociales aux méthodes computationnelles. Mes interventions portent sur :
\begin{itemize}[nosep]
    \item Programmation en Python pour l'analyse de données
    \item Analyse de réseaux sociaux avec \texttt{networkx} et \texttt{gephi}
    \item Introduction au traitement automatique des langues pour l'analyse sémantique (prétraitement, plongement lexical, transformers, classification supervisée ou non-supervisée)
\end{itemize}

\subsubsection{Analyse de réseaux pour historiens, philosophes et sociologues des sciences (Wuppertal, 2025)}

\noindent\textbf{Public :} Doctorants du Research Training Group GRK 2696\\
\textbf{Format :} Tutoriel de 2 heures

Introduction à l'analyse de réseaux appliquée à l'étude des sciences : construction de réseaux de co-autorat et de citations, détection de communautés, mesures de centralité, visualisation avec Gephi.

\subsubsection{Problèmes inverses pour philosophes (Ruhr University Bochum, 2025)}

\textbf{Responsable du cours :} Dunja Šešelja (Professeure)\\
\textbf{Format :} Cours invité de 2 heures

Intervention sur le thème \textit{``Inverse problems for philosophers: Bridging the gap between agent-based models and behavioral data''}. Présentation des méthodes d'inférence permettant de relier modèles théoriques (simulation par agents) et données comportementales observées, avec applications en épistémologie sociale.

\subsubsection{Tutoriel logiciel LongFoRMer/ChildProject (UCLA, 2024)}

\noindent\textbf{Public :} Chercheurs et doctorants en linguistique\\
\textbf{Format :} Tutoriel pratique de 2 heures

Formation à l'utilisation du package Python \texttt{ChildProject} que j'ai développé pour la gestion et l'analyse d'enregistrements audio de longue durée pour l'étude de l'acquisition de language.

\subsection{Matériel pédagogique développé}

\begin{itemize}
    \item \textbf{Documentation ChildProject} : Documentation complète du package Python, incluant tutoriels et exemples reproductibles (\url{https://childproject.readthedocs.io})
    \item \textbf{Site pédagogique cosmology.education} (2015) : Site web interactif sur l'histoire de la cosmologie moderne, avec simulations en C++ ($\sim$7\,500 utilisateurs/an) -- \url{https://cosmology.education}
    \item \textbf{Logiciel Orbit} (2012) : Programme C++ pour la résolution de problèmes à N-corps, la visualisation du mouvement planétaire et la prédiction d'éclipses -- \url{https://orbit.sciencestechniques.fr}
\end{itemize}

% \subsection{Projet d'enseignement}

% Fort de ma double formation en informatique/physique et en sciences humaines, ainsi que de mon expérience en recherche interdisciplinaire, je suis en mesure d'assurer des enseignements dans les domaines suivants :

% \textbf{Enseignements fondamentaux en informatique :}
% \begin{itemize}[nosep]
%     \item Algorithmique et programmation (Python, niveaux L1 à M2)
%     \item Structures de données et complexité algorithmique
%     \item Bases de données et SQL
%     \item Programmation orientée objet (Python, C++)
% \end{itemize}

% \textbf{Enseignements spécialisés :}
% \begin{itemize}[nosep]
%     \item Apprentissage automatique et apprentissage profond
%     \item Traitement automatique des langues (NLP)
%     \item Analyse de réseaux et graphes
%     \item Statistiques et inférence bayésienne pour l'informatique
%     \item Fouille de données et text mining
%     \item Sciences sociales computationnelles
% \end{itemize}

% Je suis particulièrement motivé par les enseignements à l'interface entre informatique et sciences humaines/sociales, répondant à une demande croissante de formation aux méthodes computationnelles dans ces disciplines.

%===============================================================================
% 5. ACTIVITÉS DE RECHERCHE
%===============================================================================

\section{Activités de recherche}

\subsection{Présentation générale}

Mes recherches portent sur la \textbf{cognition collective} : comment les collectifs génèrent, transmettent et maintiennent des connaissances complexes. J'étudie en particulier les dilemmes sociaux auxquels font face les individus dans des contextes de cognition collective, en particulier en science. Pour répondre à ces questions, je développe et applique des \textbf{méthodes computationnelles} à des données comportementales observationnelles et expérimentales.

Mon travail se situe à l'intersection de l'\textbf{informatique} (apprentissage automatique, NLP, analyse de réseaux, modélisation statistique), des \textbf{sciences cognitives} et de l'\textbf{épistémologie sociale}. Il contribue aux domaines de la \textbf{science des données} et des \textbf{sciences sociales computationnelles}.

\subsection{Axes de recherche et contributions informatiques}

\subsubsection{Axe 1 : Inférence statistique}

\subsubsection{Axe 2 : Apprentissage automatique}

\subsubsection{Axe 3 : Modélisation}




%===============================================================================
% 6. LISTE DES PUBLICATIONS
%===============================================================================

\section{Liste des publications}

\subsection{Revues internationales avec comité de lecture}

\begin{enumerate}
    \item \textbf{Gautheron, L.} (2025). Balancing specialization and adaptation in a transforming scientific landscape. \textit{EPJ Data Science}, 14, 2. \url{https://doi.org/10.1140/epjds/s13688-024-00516-8}\\
    \textit{Revue Q1 en Science des données (Scimago). Seul auteur.}
    
    \item Cristia, A., \textbf{Gautheron, L.}, Zhang, Z., Schuller, B., Scaff, C., et al. (2024). Establishing the reliability of metrics extracted from long-form recordings using LENA and the ACLEW pipeline. \textit{Behavior Research Methods}, 56, 7489--7508. \url{https://doi.org/10.3758/s13428-024-02493-2}\\
    \textit{Revue de référence en méthodologie. Contribution : analyse statistique, développement logiciel.}
    
    \item \textbf{Gautheron, L.}, Omodei, E. (2023). How research programs come apart: The example of supersymmetry and the disunity of physics. \textit{Quantitative Science Studies}, 4(3), 671--699. \url{https://doi.org/10.1162/qss_a_00262}\\
    \textit{MIT Press. Premier auteur. Contribution : conception, analyse, rédaction.}
    
    \item Cristia, A., \textbf{Gautheron, L.}, Colleran, H. (2023). Vocal input and output among infants in a multilingual context: Evidence from long-form recordings in Vanuatu. \textit{Developmental Science}, 26(3), e13375. \url{https://doi.org/10.1111/desc.13375}\\
    \textit{Contribution : analyse statistique, développement des outils d'analyse.}
    
    \item \textbf{Gautheron, L.}, Rochat, N., Cristia, A. (2022). Managing, storing, and sharing long-form recordings and their annotations. \textit{Language Resources and Evaluation}, 56, 1--27. \url{https://doi.org/10.1007/s10579-022-09579-3}\\
    \textit{Springer. Premier auteur. Article de présentation du package ChildProject.}
    
    \item Lavechin, M., Seyssel, M., \textbf{Gautheron, L.}, Dupoux, E., Cristia, A. (2022). Reverse Engineering Language Acquisition with Child-Centered Long-Form Recordings. \textit{Annual Review of Linguistics}, 8, 389--407. \url{https://doi.org/10.1146/annurev-linguistics-031120-122120}\\
    \textit{Article de revue dans Annual Reviews.}
\end{enumerate}

\subsection{Prépublications (preprints)}

\begin{enumerate}
    \item \textbf{Gautheron, L.}, Kidd, E., Malko, A., Lavechin, M., Cristia, A. (2025). Classification errors distort findings in automated speech processing: examples and solutions from child-development research. \textit{arXiv:2508.15637}. \url{https://arxiv.org/abs/2508.15637}\\
    \textit{Premier auteur. Soumis à revue internationale.}
    
    \item \textbf{Gautheron, L.} (2025). Dilemmas and trade-offs in the diffusion of conventions. \textit{arXiv:2501.17300}. \url{https://arxiv.org/abs/2501.17300}\\
    \textit{Seul auteur. Soumis à revue internationale.}
\end{enumerate}

\subsection{Conférences internationales avec comité de lecture}

\begin{enumerate}
    \item \textbf{Gautheron, L.} (2025). Dilemmas and trade-offs in the diffusion of conventions. \textit{11th International Conference on Computational Social Science (IC2S2)}, Norrköping, Suède. [Présentation orale; taux d'acceptation: 25\%]
    
    \item \textbf{Gautheron, L.}, Schneider, M.D. (2025). Plural pursuit across scales. \textit{Workshop on Large Language Models for the History, Philosophy, and Sociology of Science}, TU Berlin. [Présentation orale]
    
    \item \textbf{Gautheron, L.} (2024). Balancing Specialization and Adaptation in a Transforming Scientific Landscape. \textit{10th International Conference on Computational Social Science (IC2S2)}, Philadelphia, États-Unis. [Poster]
    
    \item \textbf{Gautheron, L.} (2022). The many faces of supersymmetry: Supersymmetry across subcultures of High-Energy Physics, 1971--2019. \textit{History of Science Society Annual Meeting}, Chicago, États-Unis. [Présentation orale]
\end{enumerate}

\subsection{Conférences invitées}

\begin{enumerate}
    \item \textbf{Gautheron, L.} (2025). A statistical physics approach to dilemmas and trade-offs challenging the diffusion of conventions. \textit{Institut Jean-Nicod}, ENS.
    
    \item \textbf{Gautheron, L.} (2024). ``When her family finds you are using the wrong metric...': dilemmas and trade-offs in the diffusion of scientific conventions. \textit{Social Epistemology Research Group}, London School of Economics.
    
    \item \textbf{Gautheron, L.} (2024). Correlational analyses of the effects of caregivers' speech behaviour on children's speech production. \textit{Department of Linguistics}, UCLA.
    
    \item \textbf{Gautheron, L.} (2023). Balancing Specialization and Adaptation in a Transforming Scientific Landscape. \textit{NLP Seminar}, LATTICE, Montrouge.
    
    \item \textbf{Gautheron, L.} (2023). Too beautiful to be false, or too beautiful to be true: supersymmetry and the future of high-energy physics. \textit{2JM Seminar}, Sciences Po Paris.
\end{enumerate}

\subsection{Productions logicielles}

\begin{enumerate}
    \item \textbf{ChildProject} (2020--présent) : Package Python pour la gestion et l'analyse d'enregistrements audio longue durée. \url{https://github.com/LAAC-LSCP/ChildProject}\\
    \textit{Développeur principal. 15 étoiles GitHub, documentation complète, tests automatisés.}
    
    \item \textbf{diarization-simulation} (2025) : Package Python pour la simulation d'algorithmes de diarisation. \url{https://github.com/LAAC-LSCP/diarization-simulation}\\
    \textit{Développeur unique.}
    
    \item \textbf{Code reproductible} : Tous les codes et données associés à mes publications sont disponibles sur GitHub avec documentation.
\end{enumerate}

\subsection{Travaux joints au dossier (3 publications les plus significatives)}

\noindent\textbf{[1] Gautheron, L. (2025). Balancing specialization and adaptation in a transforming scientific landscape. \textit{EPJ Data Science}.}

\textit{Contribution propre en informatique :} Article en auteur unique. Développement d'un modèle bayésien original combinant transport optimal inverse et topic modeling pour analyser l'évolution des intérêts de recherche de 2\,108 scientifiques sur 20 ans. Contributions techniques : (1) modèle probabiliste de transport optimal pour l'adaptation des portfolios de recherche, (2) inférence MCMC des coûts d'apprentissage avec Stan, (3) traitement NLP de 186\,162 résumés scientifiques, (4) analyse de réseaux de co-autorat pour mesurer le capital social.

\vspace{0.3cm}

\noindent\textbf{[2] Gautheron, L., Omodei, E. (2023). How research programs come apart. \textit{Quantitative Science Studies}.}

\textit{Contribution propre en informatique :} Premier auteur (80\% du travail). Développement de méthodes computationnelles pour l'analyse de la fragmentation scientifique : (1) pipeline d'extraction et de nettoyage de données bibliographiques, (2) application de topic models (LDA) pour identifier les sous-communautés, (3) construction et analyse de réseaux de citations et co-autorat, (4) mesures quantitatives de diversité et de cohésion.

\vspace{0.3cm}

\noindent\textbf{[3] Gautheron, L., Rochat, N., Cristia, A. (2022). Managing, storing, and sharing long-form recordings. \textit{Language Resources and Evaluation}.}

\textit{Contribution propre en informatique :} Premier auteur (70\% du travail). Conception et développement complet du package Python ChildProject : (1) architecture de gestion de données pour corpus audio massifs, (2) système de validation et contrôle qualité, (3) intégration de pipelines de diarisation automatique, (4) API et interface en ligne de commande, (5) documentation et tests automatisés.

%===============================================================================
% 7. ENCADREMENTS
%===============================================================================

%===============================================================================
% 8. RESPONSABILITÉS COLLECTIVES ET ADMINISTRATIVES
%===============================================================================

\section{Responsabilités collectives et administratives}

\subsection{Organisation d'événements scientifiques}

\begin{itemize}
    \item \textbf{2024} : Co-organisateur du workshop \textit{``Methodological Transformations in Fundamental Physics''} (16--18 septembre 2024, Wuppertal, Allemagne)
\end{itemize}

\subsection{Activités d'évaluation}

\begin{itemize}
    \item \textbf{Relecture pour revues :} \textit{Behavior Research Methods}
    \item \textbf{Relecture pour conférences :} \textit{International Conference on Computational Social Science} (IC2S2)
\end{itemize}

\subsection{Participation à des projets de recherche}

\begin{itemize}
    \item \textbf{2025--2027} : Membre du projet NSF \textit{``Designing Smart Environments to Augment Collective Learning \& Creativity''} (PIs : D. Conley (Princeton), S. Frey (UC Davis), N. Jacoby (Cornell), O. Tchernichovski (CUNY))
\end{itemize}

\subsection{Financements obtenus}

\begin{itemize}
    \item \textbf{2025} : Co-investigateur, \textit{``CASCade: Building a Center for Applications in Social Complexity''}, University of Missouri, ``Big-ideas'' grant (\$53\,000)
    \item \textbf{2023} : G-Research Grant for PhD students and postdocs in quantitative fields (£1\,000)
\end{itemize}

%===============================================================================
% 9. FORMATIONS ET ÉCOLES D'ÉTÉ
%===============================================================================

\section{Formations et écoles d'été}

\begin{itemize}
    \item \textbf{2025} : Santa Fe Institute \textit{Advanced} Graduate Workshop in Computational Social Science Modeling and Complexity, Santa Fe, NM, États-Unis
    \item \textbf{2024} : Santa Fe Institute Graduate Workshop in Computational Social Science Modeling and Complexity, Santa Fe, NM, États-Unis (taux d'acceptation : 20\%)
    \item \textbf{2023} : École d'été \textit{``Collaboration and Interdisciplinarity in Science and Technology''}, Wuppertal
    \item \textbf{2022} : École de printemps \textit{``The History, Philosophy and Sociology of Large Scale Experiments''}, 4\textsuperscript{e} International Spring School of the Epistemology of the Large Hadron Collider, Wuppertal
    % \item \textbf{2021} : Formation CNRS \textit{``Basics of Machine Learning and Deep Learning''} (28h)
\end{itemize}

%===============================================================================
% 10. SÉJOURS DE RECHERCHE
%===============================================================================

\section{Séjours de recherche}

\begin{itemize}
    \item \textbf{Octobre 2025} : Cornell University, Ithaca, États-Unis
    \item \textbf{Février--Août 2023} : Medialab, Sciences Po Paris (étudiant visiteur)
    \item \textbf{Juin 2022} : Max Planck Institute for Evolutionary Anthropology, Leipzig, Allemagne
\end{itemize}

%===============================================================================
% 11. COMPÉTENCES TECHNIQUES
%===============================================================================

\section{Compétences techniques}

\subsection{Programmation}

\begin{itemize}
    \item \textbf{Expert :} Python (numpy, pandas, scikit-learn, pytorch, networkx, graph-tool, gensim, bertopic, pyannote)
    \item \textbf{Avancé :} C, C++, Fortran, Stan (inférence bayésienne), SQL
    \item \textbf{Intermédiaire :} JavaScript, PHP, HTML/CSS, Arduino
    \item \textbf{Outils :} Git, GitHub Actions, Docker, Linux, LaTeX
\end{itemize}

\subsection{Méthodes}

\begin{itemize}
    \item Apprentissage automatique supervisé et non supervisé
    \item Inférence bayésienne (MCMC, HMC, variationnel)
    \item Traitement automatique des langues (topic models, transformers, sentence embeddings)
    \item Analyse de réseaux (détection de communautés, stochastic block models)
    \item Transport optimal et problèmes inverses
    \item Traitement du signal audio
\end{itemize}

\subsection{Langues}

\begin{itemize}
    \item Français : langue maternelle
    \item Anglais : courant (C2)
    \item Espagnol : notions
\end{itemize}

%===============================================================================
% 12. RÉFÉRENCES
%===============================================================================

\section{Références}

\begin{tabular}{@{}p{0.48\textwidth}p{0.48\textwidth}@{}}

\textbf{Alejandrina Cristia} & \textbf{Elisa Omodei} \\
Directrice de recherche CNRS & Professeure \\
LSCP, Département d'Études Cognitives & Department of Network and Data Science \\
École Normale Supérieure, Paris & Central European University, Vienne \\
\href{mailto:alecristia@gmail.com}{alecristia@gmail.com} & \href{mailto:omodeie@ceu.edu}{omodeie@ceu.edu} \\
\\
\textbf{Thomas Heinze} & \textbf{Radin Dardashti} \\
Professeur (Directeur de thèse) & Professeur (Directeur de thèse) \\
IZWT, Bergische Universität Wuppertal & IZWT, Bergische Universität Wuppertal \\
\href{mailto:theinze@uni-wuppertal.de}{theinze@uni-wuppertal.de} & \href{mailto:radin@gmx.de}{radin@gmx.de} \\
\\
\textbf{Olivier Darrigol} & \\
Directeur de recherche émérite CNRS & \\
SPHERE, Université Paris-Cité & \\
\href{mailto:darrigol@paris7.jussieu.fr}{darrigol@paris7.jussieu.fr} & \\

\end{tabular}

\end{document}
